% Nouveau EGA V (cote 25): archive cote 25, site page 105 (PDF page 106).
% typescript. Draft compared: none.
% AI-transcribed by Claude Opus 5.5 (claude-opus-5-5, adaptive thinking, effort high), Message Batches API, batch ega25b (msgbatch_01A7n5YuXubUPNK1zsepi8RG), 2026-11-13; checked against both supplied scan parts; see CHECK_DECISIONS.tsv.
\scanpage{25/105}
\mspage{8}

de $\Gamma(U,\underline{O}_U)^{\hand{*}}$ .

En particulier:

\uline{Corollaire} 1.8. Soient S un préschéma, P=P(\underline{E}) un fibré projectif sur S
défini par \underline{E} loc libre de type fini, u un automorphisme de P. Alors u
est déterminé, au vosinage de tout point s $\hand{\in}$ S , par un automorphisme
u de \underline{E}\hand{|}U, ce dernier étant bien déterminé par u modulo multiplication par
un élément de $\Gamma(U,\underline{O}_U)^{\hand{*}}$ .

% NOTE: in the next line the words "on déduit peut déduire facilement" are circled by hand
\uline{Remarque} 1.9. De 1.8. on  peut déduire facilement que le foncteur \placedmargin{\uline{le faire} ! \unsure{C|c}'est l'ouvert de \unsure{$P(\underline{E}^{\vee}\otimes\underline{E})$|$P(\underline{E}\otimes\underline{E})$} défini \unsure{pr|par} la non annulation du ``déterminant'' - - - -}
 en groupes $\underline{Aut}_S(P)$ sur S est représentable par un p\unsure{r|e}éschéma affine
de présentation finie sur S, qui s'interprète aussi comme le préschéma en
groupes \uline{quotient} du schéma en groupes linéaire $\underline{Gl}(\underline{E})$ par son centre
$\underline{G}_m$ . Ce préschéma en groupes s'appelle 
le \uline{préschéma en groupes projectif}, ou simplement \uline{groupe projectif}, défini
par \underline{E}, et se note GP(\underline{E})\hand{,} Lorsque \underline{E} est libre, $\underline{E}\hand{\simeq}\underline{O}_S^{r+1}$, alors GP(\underline{E} ) n'est autre que le
 préschéma en groupes $GP(r)_S$ déduit par le changement de base
$S\hand{\to}\mathrm{Spec}(\underline{Z})$  du schéma en groupes analogue GP(r) sur Spec(\underline{Z}), appelé
le \uline{Groupe \unsure{P|p}rojectif} \uuline{absolu}.

\placedmargin{105}
